\documentclass{article}
\usepackage[utf8]{inputenc}
\usepackage{a4wide}

\begin{document}

\section*{เกมบันไดงู}

กำหนดจำนวนช่องของกระดาน เช่น100ให้หาจำนวนครั้ง{\textbf{\ul{ขั้นต่ำ}}}ที่ต้องเล่น

(ทอยลูกเต๋า 6 หน้า) เพื่อให้จบเกม โดยกำหนดบันได (ladder) เพื่อก้าวไปข้างหน้า และงู

(snake) ที่ทำให้ตกถอยหลัง ไว้ในโปรแกรม~

ทั้งนี้เริ่มต้นการเล่นจะอยู่ช่องที่\texttt{1}และเกมจบเมื่อไปถึงช่องที่\texttt{100}



\includegraphics[width=2.08333in,height=2.08333in,alt={Screenshot 2567-03-28 at 11.34.55.png}]{/public/upload/beb5b8f3d5.png}\\



\textbf{ตัวอย่าง}



{\def\LTcaptype{none} % do not increment counter

\begin{longtable}[]{@{}

  >{\raggedright\arraybackslash}p{(\linewidth - 4\tabcolsep) * \real{0.3333}}

  >{\raggedright\arraybackslash}p{(\linewidth - 4\tabcolsep) * \real{0.3333}}

  >{\raggedright\arraybackslash}p{(\linewidth - 4\tabcolsep) * \real{0.3333}}@{}}

\toprule\noalign{}

\begin{minipage}[b]{\linewidth}\raggedright

ข้อมูลเข้า

\end{minipage} & \begin{minipage}[b]{\linewidth}\raggedright

ข้อมูลออก

\end{minipage} & \begin{minipage}[b]{\linewidth}\raggedright

คำอธิบาย

\end{minipage} \\

\midrule\noalign{}

\endhead

\bottomrule\noalign{}

\endlastfoot

\begin{minipage}[t]{\linewidth}\raggedright

100 8 8\\

1 38\\

4 14\\

9 31\\

21 42\\

28 84\\

51 67\\

72 91\\

80 99\\

17 7\\

54 34\\

62 19\\

64 60\\

87 36\\

93 73\\

95 75\\

98 79\\

\strut

\end{minipage} & 7 & \begin{minipage}[t]{\linewidth}\raggedright

กระดานมีขนาด 100 ช่อง (10 x 10) มีบันได 8 ชุด และงู 8 ตัว บันไดชุดแรกอยู่ที่ 1

เพื่อกระโดดไป 38 ทันที ส่วนงูมีอยู่ที่ช่อง\textbf{17 54 62 64 87 93

95}และ\textbf{98}งูตัวสุดท้ายจะทำให้ตกลงไปที่ช่อง\textbf{79}ทันที\\

\strut

\end{minipage} \\

\begin{minipage}[t]{\linewidth}\raggedright

100 ~0 0\\

\strut

\end{minipage} & \begin{minipage}[t]{\linewidth}\raggedright

17\\

\strut

\end{minipage} & \begin{minipage}[t]{\linewidth}\raggedright

กระดานมีขนาด 100 ช่อง (10 x 10) ไม่มีบันไดและงูเลย\\

\strut

\end{minipage} \\

\end{longtable}

}



\subsection*{Input}
มีจำนวน\texttt{1\ +\ P\ +\ Q}บรรทัด ดังนี้\\



{\def\LTcaptype{none} % do not increment counter

\begin{longtable}[]{@{}

  >{\raggedright\arraybackslash}p{(\linewidth - 2\tabcolsep) * \real{0.5000}}

  >{\raggedright\arraybackslash}p{(\linewidth - 2\tabcolsep) * \real{0.5000}}@{}}

\toprule\noalign{}

\begin{minipage}[b]{\linewidth}\raggedright

บรรทัดที่ 1\\

\strut

\end{minipage} & \begin{minipage}[b]{\linewidth}\raggedright

จำนวนเต็ม 3 จำนวนคั่นด้วยช่องว่างค่าแรกเป็นจำนวนช่องของเกม N ค่าที่สองเป็นค่า P

คือจำนวนบันไดและค่าที่สามเป็นค่าQ คือจำนวนงู~ , 0 \textless= P,Q \textless= 5,000

; 10 \textless= N \textless= 10,000\\

\strut

\end{minipage} \\

\midrule\noalign{}

\endhead

\bottomrule\noalign{}

\endlastfoot

\begin{minipage}[t]{\linewidth}\raggedright

\textbf{บรรทัดที่ 2 ถึง 1+P\\

}\strut

\end{minipage} & \begin{minipage}[t]{\linewidth}\raggedright

บอกบันไดจำนวน\textbf{P}บรรทัด แต่ละบรรทัดเป็นจำนวนเต็ม 2 จำนวน คั่นด้วยช่องว่าง

ค่าแรกบอกช่องต้นบันได ค่าที่สองบอกช่องปลายบันได โดยที่ช่องปลายบันไดจะมากกว่าช่องต้นบันได

หากเดินมาตกที่ช่องต้นของบันได ก็จะข้ามไปช่องปลายบันไดทันที ,\textbf{0 \textless=

Pi\textless{} N}\\

\strut

\end{minipage} \\

\begin{minipage}[t]{\linewidth}\raggedright

\textbf{บรรทัดที่ 2+P ถึง 1+P+Q\\

}\strut

\end{minipage} & \begin{minipage}[t]{\linewidth}\raggedright

บอกงูจำนวน Q บรรทัด แต่ละบรรทัดเป็นจำนวนเต็ม 2 จำนวน คั่นด้วยช่องว่าง\\

ค่าแรกบอกช่องหัวงู ค่าที่สองบอกช่องหางงู โดยที่ช่องหางงูจะน้อยกว่าช่องหัวงู

หากเดินตกที่ช่องหัวงู ก็จะตกไปช่องหางงูทันที ,\textbf{0 \textless= Qi\textless{}

N}\\

\strut

\end{minipage} \\

\end{longtable}

}



\subsection*{Output}
จำนวนครั้ง{ขั้นต่ำ}ที่ต้องเล่น เป็นเลขจำนวนเต็ม\\



\end{document}
